\documentclass[11pt]{article}
\usepackage[margin=1in]{geometry}
\usepackage{amsmath,amssymb}
\usepackage{parskip}
\usepackage[T1]{fontenc}
\usepackage{enumitem}
\usepackage{titlesec}

\titleformat{\section}{\normalsize\bfseries}{}{0pt}{}
\titlespacing{\section}{0pt}{12pt}{4pt}
\setlist[itemize]{leftmargin=1.4em,topsep=2pt,itemsep=3pt}
\newcommand{\rf}[1]{\hfill\textnormal{\small[#1]}}

\begin{document}

\begin{center}
{\Large\bfseries List of Changes --- Manuscript XQ11175}\\[4pt]
\textit{``How Focused Are LLMs? A Quantitative Study via Repetitive Deterministic Prediction Tasks''}
\end{center}

\medskip

All changes respond to specific referee comments (R1 = first referee, R2 = second referee). No
content was added beyond what the referees requested.

\section{Abstract}
\begin{itemize}
  \item Softened the description of the statistical model to an ``effective'' description linking
        correlated token errors to sequence-level failure (removed the implied first-principles
        attention mechanism). \rf{R2-1}
\end{itemize}

\section{Introduction}
\begin{itemize}
  \item Rewrote the ``factorizable / product of per-step accuracies'' statement to be precise:
        the factorization is exact only for site-local tasks (letter replacement, Pauli
        multiplication); integer addition has long-range carry dependence and is only
        approximately factorizable. \rf{R1-1}
  \item Added a sentence framing the spin-glass analogy as an effective, phenomenological
        description, not a microscopic theory of the network. \rf{R1-3, R2-1}
\end{itemize}

\section{Experiment (Methods)}
\begin{itemize}
  \item Stated the sample size: up to $n=100$ independent instances per length (10 repetitions
        $\times$ batch 10); noted that a subset of grok-4 runs use fewer repetitions
        ($n\sim20$--$30$). \rf{R1-2, R2-3}
  \item Stated that all accuracy curves report SAR with $95\%$ confidence-interval error bars:
        bootstrap over the repeated batches for the geometric-mean SAR, and Wilson score
        intervals where SAR is the instance-level success rate (integer addition and the scaling
        collapse). \rf{R1-2, R2-3}
  \item Strengthened the no-tools statement: all queries used plain text completion/chat calls
        with no tools, function-calling, or code-interpreter enabled; results reflect in-context
        computation only. \rf{R2-4}
  \item Stated the decoding settings (temperature 0 where exposed; reasoning models use fixed
        internal sampling; per-model settings in the released code). \rf{R1/R2 general}
\end{itemize}

\section{Figures}
\begin{itemize}
  \item All accuracy-vs-length figures (cyclic letter replacement, integer addition, Pauli
        multiplication) now display $95\%$ confidence intervals (bootstrap for the geometric-mean
        SAR curves; Wilson for the instance-level success rate, i.e.\ integer addition and the
        scaling collapse) and state $n$ in the captions. \rf{R1-2, R2-3}
  \item Added a new scaling-collapse figure: SAR vs $N/N_*$ collapses onto a single crossover
        curve across all models and tasks, confirming the universal finite-size crossover.
        \rf{R2-2}
  \item Phase (correlation-error) maps regenerated from the updated fits.
\end{itemize}

\section{Analysis (Empirical formula)}
\begin{itemize}
  \item Added a model-comparison statement: an AIC comparison against the single-exponential
        ($\alpha=1$) baseline favors the double-exponential in all but the highest-accuracy
        cases, which sit in the $\alpha\to1$ limit; the per-curve comparison is in the released
        code. \rf{R2-5}
\end{itemize}

\section{Analysis (Hypothesis / statistical model)}
\begin{itemize}
  \item Reframed the couplings $J_{ij}$ as an effective form inspired by (not derived from)
        attention; removed causal ``attention-induced/attention-mediated'' language. \rf{R2-1}
  \item Added an argument that disordered couplings are necessary: a field-only model (uniform or
        random $h_i$, no coupling) factorizes $\rightarrow$ single exponential, no crossover;
        uniform ferromagnetic couplings give a smooth transition, not the observed sharp,
        sample-to-sample crossover. \rf{R1-4}
  \item Added clarification on symmetry and energetics: the field $h$ breaks the $s\to-s$
        symmetry (no $\mathbb{Z}_2$ degeneracy); at large $N$ the disorder energy ($\sim N^2$)
        overtakes the field, so an incorrect configuration can become energetically competitive;
        the frustration is the competition between random-sign couplings and the ordering field.
        \rf{R1-3}
  \item Added a sentence introducing the scaling-collapse prediction and referencing the new
        figure. \rf{R2-2}
  \item Listed microscopic connection (extract $J_{ij}$ from attention; perturb attention
        sparsity) as future work. \rf{R2-1}
\end{itemize}

\section{Divide-and-conquer section}
\begin{itemize}
  \item Stated that the divide-and-conquer prediction follows from the empirical accumulation law
        and does not depend on the SK-specific equations. \rf{R1-5a}
  \item Attributed the Pro/Flash difference to per-call overhead and no longer present the Flash
        curve as a success. \rf{R1-5b}
  \item Flagged extension of the divide-and-conquer test to the other tasks as future work.
        \rf{R1-5c}
\end{itemize}

\section{Appendix A (SK derivation)}
\begin{itemize}
  \item Corrected the $\mathcal{O}(J_0^4)$ coefficient in the large-field reduction from
        $+2J_0^4(N-1)^2$ to $-4J_0^4(N-1)^2$. \rf{R1-6, R2-6}
  \item Reconciled the two functional forms: the double-exponential is the exponentiation of the
        leading-order $\mathcal{O}(J_0^2)$ correction (a phenomenological closed form valid for
        small $J_0$, large $h$), NOT a strict term-by-term resummation; the leading-order
        $(\alpha,\beta_0)\leftrightarrow(J_0,h)$ mapping is retained. \rf{R1-6, R2-6}
  \item Fixed a broken equation reference (``Eq.~(0)'') to the explicit token-accuracy relation.
\end{itemize}

\section{Conclusion / Future work}
\begin{itemize}
  \item Added two future directions tied to referee suggestions: (i)~extending divide-and-conquer
        to the other tasks; (ii)~connecting the model to network internals (extracting $J_{ij}$
        from attention; testing attention-sparsity perturbations). \rf{R1-5c, R2-1}
\end{itemize}

\section{Notes}
\begin{itemize}
  \item The manuscript and response use gender-neutral terms throughout.
  \item No new experiments were run; all strengthening derives from re-analysis of the
        already-collected data, code-level clarification, and revised framing.
\end{itemize}

\end{document}
